\begin{mistake}{2026-04-12}{10}

\begin{problemcontent}
设
\[
A=[\alpha_1,\alpha_2,\alpha_3,\alpha_4]
\]
为四阶矩阵，经初等行变换变为
\[
B=\begin{bmatrix}
1&0&-1&2\\
0&1&2&3\\
0&0&0&0\\
0&0&0&0
\end{bmatrix},
\]
求满足
\[
A=(\alpha_1,\alpha_2)K
\]
的矩阵 \(K\)。

\end{problemcontent}

\begin{solutioncontent}
由于 \(A\) 经过的是初等行变换，列向量之间的线性关系保持不变。因此在 \(B\) 中各列如何由前两列线性表示，在 \(A\) 中也完全相同。

观察 \(B\) 的列向量：
\[
\beta_1=\begin{pmatrix}1\\0\\0\\0\end{pmatrix},
\quad
\beta_2=\begin{pmatrix}0\\1\\0\\0\end{pmatrix},
\quad
\beta_3=\begin{pmatrix}-1\\2\\0\\0\end{pmatrix},
\quad
\beta_4=\begin{pmatrix}2\\3\\0\\0\end{pmatrix}.
\]

于是有
\[
\beta_3=-\beta_1+2\beta_2,
\qquad
\beta_4=2\beta_1+3\beta_2.
\]

故在原矩阵中同样成立
\[
\alpha_3=-\alpha_1+2\alpha_2,
\qquad
\alpha_4=2\alpha_1+3\alpha_2.
\]

而显然
\[
\alpha_1=1\alpha_1+0\alpha_2,
\qquad
\alpha_2=0\alpha_1+1\alpha_2.
\]

因此
\[
A=(\alpha_1,\alpha_2)
\begin{bmatrix}
1&0&-1&2\\
0&1&2&3
\end{bmatrix}.
\]

所以
\[
K=\begin{bmatrix}
1&0&-1&2\\
0&1&2&3
\end{bmatrix}.
\]

\end{solutioncontent}

\mistakeknowledge{
\begin{itemize}
  \item \textbf{核心公式/定理}：初等行变换不改变列向量之间的线性关系，因此行最简形中非主元列对主元列的表示系数，可直接迁移回原矩阵。
  \item \textbf{易错点}：只有“行变换”保持列关系，“列变换”保持行关系，二者不能混用。另需注意矩阵维数：\((\alpha_1,\alpha_2)\) 为 \(4\times 2\)，故 \(K\) 必为 \(2\times 4\)。
  \item \textbf{关联知识}：行最简形中主元列对应极大线性无关组，非主元列坐标正是表示系数，这一结论在线性表示与秩的题中非常高频。
\end{itemize}}

\end{mistake}
